% =====================================================================
%  preambulo.tex — configuración común
%  Compartido por el documento maestro y por los envoltorios de sección
%  (sec-01… sec-04). Se incluye siempre desde documento/, de modo que
%  las rutas relativas apuntan a ../diagramas/.
% =====================================================================
\documentclass[11pt,spanish]{article}

\usepackage[utf8]{inputenc}
\usepackage[T1]{fontenc}
\usepackage{lmodern}
\usepackage[spanish,es-nodecimaldot,es-tabla]{babel}
\usepackage[a4paper,margin=2.5cm]{geometry}
\usepackage{graphicx}
\usepackage{booktabs}
\usepackage{longtable}
\usepackage{array}
\usepackage{float}
\usepackage{placeins}
\usepackage{hyperref}
\usepackage{csquotes}
\usepackage[table]{xcolor}
\usepackage{pdflscape}

\definecolor{FichaHeaderBg}{HTML}{2F3B47}
\definecolor{FichaBorder}{HTML}{94A3B8}

% Semáforo del registro de riesgos (sección 6)
\definecolor{EstadoResuelto}{HTML}{D8EBDC}
\definecolor{EstadoEnProceso}{HTML}{FFF6DB}
\definecolor{EstadoPendiente}{HTML}{EEF0F2}
\definecolor{EstadoRetrasado}{HTML}{F7D5D3}

\graphicspath{{../diagramas/}{../interfaces/}}
\renewcommand{\arraystretch}{1.25}

% Compilacion reproducible: el PDF versionado no cambia por fecha ni por
% metadatos de ruta, de modo que dos compilaciones producen el mismo archivo.
\pdfinfoomitdate=1
\pdftrailerid{}
\pdfsuppressptexinfo=-1

\hypersetup{
  pdftitle={Requisitos del Proyecto RedPepper},
  pdfauthor={Gabriel Pajuelo, Armando Enriquez, Victor Lupinta, Jose Chino, Jose Choquehuanca},
  pdfsubject={Ingeniería de Requisitos},
  colorlinks=true,
  linkcolor=black,
  urlcolor=blue,
  citecolor=black
}

% Columnas reutilizables
\newcolumntype{L}[1]{>{\raggedright\arraybackslash}p{#1}}
\newcolumntype{R}[1]{>{\raggedleft\arraybackslash}p{#1}}

% Entorno de ficha (Campo | Contenido) segun plantillas de Ingenieria de Requisitos
\newenvironment{ficha}[1]{%
  \par\vspace{0.45em}\noindent
  \begin{minipage}{\linewidth}
  \textbf{#1}\par\vspace{0.25em}
  \small
  \arrayrulecolor{FichaBorder}%
  \begin{tabular}{|L{3.9cm}|L{11.1cm}|}
  \hline
  \rowcolor{FichaHeaderBg}\textbf{\color{white}Campo} & \textbf{\color{white}Contenido} \\
  \hline
}{%
  \end{tabular}%
  \arrayrulecolor{black}%
  \end{minipage}\par\vspace{0.55em}
}
\newcommand{\fichafila}[2]{\textbf{#1} & #2 \\ \hline}
